\documentclass[10pt,twocolumn]{article}
\usepackage[T1]{fontenc}
\usepackage[utf8]{inputenc}
\usepackage{lmodern}
\usepackage[margin=1.8cm,top=2.4cm,bottom=2cm,columnsep=0.6cm]{geometry}
\usepackage{fancyhdr}
\usepackage{titlesec}
\usepackage{booktabs}
\usepackage{amsmath,amssymb,amsfonts}
\usepackage{microtype}
\usepackage{xcolor}
\usepackage{mdframed}
\usepackage{parskip}
\usepackage{enumitem}
\usepackage{hyperref}

\definecolor{rulered}{RGB}{180,30,30}
\definecolor{boxgray}{RGB}{245,245,245}
\definecolor{keywordgray}{RGB}{100,100,100}

\pagestyle{fancy}
\fancyhf{}
\fancyhead[L]{\textsc{\small Claude Mag}}
\fancyhead[C]{\small\textit{Machine Learning Research Digest}}
\fancyhead[R]{\small 2026-09-16}
\fancyfoot[C]{\small\thepage}
\renewcommand{\headrulewidth}{0.8pt}

\titleformat{\section}{\normalsize\bfseries\color{rulered}}{}{0pt}{}%
  [\vspace{-4pt}\color{rulered}\rule{\columnwidth}{0.4pt}]
\titlespacing{\section}{0pt}{8pt}{4pt}

\hypersetup{colorlinks=true,urlcolor=rulered,linkcolor=black}

\begin{document}
\setlength{\parindent}{0pt}
\setlength{\parskip}{4pt}

{\small\color{keywordgray}\textbf{Keywords:}
  recurrent neural networks \textbullet{}
  depth vs.\ width \textbullet{}
  compute allocation \textbullet{}
  sequence modeling \textbullet{}
  neural scaling}\par\vspace{4pt}

{\LARGE\bfseries\raggedright
  Temporal Recurrence Favors
  Fewer Layers}\par\vspace{4pt}

{\large\itshape\raggedright
  Recurrent Models Peak at 2--4 Layers While Non-Recurrent Models
  Keep Improving Up to 32 --- Temporal Depth and Within-Step
  Depth Provide Overlapping Computational Benefits}\par\vspace{6pt}

{\small\textbf{By} Ivan Anokhin, Johan Obando-Ceron, Irina Rish,
  Sebastian Risi
  (Mila / Universit\'e de Montr\'eal; University of Copenhagen)}\par\vspace{2pt}

{\small\color{keywordgray}arXiv:2609.12531}
\par\vspace{2pt}\noindent{\color{rulered}\rule{\columnwidth}{0.6pt}}\vspace{6pt}

Deep learning lore holds that depth is essential for model expressivity.
This intuition was developed for non-recurrent models that process all
sequence positions in parallel and must instantiate all computation
simultaneously. Recurrent architectures add a qualitatively different
computation axis: temporal recurrence, which accumulates processing
across sequence positions step by step. This paper asks: when temporal
recurrence is available, how much within-step depth is still needed?
The answer is surprisingly little --- and the difference has direct
implications for how recurrent models should be scaled.

\begin{mdframed}[backgroundcolor=boxgray,linecolor=rulered,linewidth=1pt,
    innerleftmargin=6pt,innerrightmargin=6pt,
    innertopmargin=6pt,innerbottommargin=6pt]
\textbf{\small AT A GLANCE}\par\vspace{4pt}
\begin{description}[leftmargin=0pt,labelsep=4pt,itemsep=2pt,parsep=0pt,
    font=\small\bfseries]
  \item[Problem:] How should compute be allocated across within-step
    depth, expert width, and parallel experts for recurrent vs.\
    non-recurrent models?
  \item[Why it matters:] Miscalibrated depth wastes compute that
    could go to width; the right allocation changes at optimal
    recurrent model design.
  \item[Finding:] Recurrent models peak at 2--4 layers; non-recurrent
    models keep improving through 8--32 layers.
  \item[Explanation:] Temporal recurrence and within-step depth both
    allow later computation to condition on earlier outputs ---
    recurrence makes depth partly redundant.
  \item[Evidence:] Sokoban planning and FineWeb language modeling;
    matched FLOPs across depth, width, expert-count axes.
\end{description}
\end{mdframed}

\section{The Big Picture}

A non-recurrent Transformer of depth $D'$ processes each token through
$D'$ layers applied once in parallel. A recurrent model of within-step
depth $D$ applied over $T$ sequence positions uses $D\cdot T$ total
layer applications. Under matched FLOPs, the question is: for a fixed
$D\cdot T$ budget, should $D$ be large (deep per-step, fewer steps of
computation) or small (shallow per-step, many recurrent steps)?

The answer the authors find empirically: small $D$ with compensating
width is consistently better for recurrent models, while large $D$ is
better for non-recurrent models. This suggests that within-step depth
and temporal depth are not fully orthogonal dimensions --- they provide
partly overlapping computational benefits.

\section{Technical Formulation}

Let $\mathbf{h}_t^{(0)} = \mathbf{x}_t$ be the input at position $t$.
The recurrent forward pass applies $D$ within-step layers:
\[
  \mathbf{h}_t^{(d)} = f_d\!\left(\mathbf{h}_t^{(d-1)},\,
  \mathbf{h}_{t-1}^{(D)}\right), \quad d = 1,\ldots,D
\]
where $\mathbf{h}_{t-1}^{(D)}$ is the hidden state from the previous
position, and weights are shared across positions (weight tying across
time). For $E$ parallel experts:
\[
  \mathbf{h}_t^{(d)} = \sum_{e=1}^E \pi_e(\mathbf{h}_t^{(d-1)})
  \cdot f_d^{(e)}\!\left(\mathbf{h}_t^{(d-1)}\right)
\]
\textbf{Compute matching:} When depth is reduced from $D$ to $D/k$,
the hidden dimension is increased by $\sqrt{k}$ to hold parameters
constant (quadratic scaling of attention/MLP with hidden dim).

\textbf{Arithmetic intensity:} Total layer applications per sequence
equal $D\cdot T$ (recurrent) vs.\ $D'$ (non-recurrent, applied once
to all positions in parallel). FLOPs are matched across configurations.

\section{What the Results Show}

\subsection*{Sokoban Planning (Success Rate, \%)}

\begin{center}
\begin{tabular}{lcc}
\toprule
Depth $D$ & Recurrent & Non-Recurrent \\
\midrule
1  & 44.2 & 31.8 \\
2  & 67.9 & 49.3 \\
4  & \textbf{71.3} & 58.7 \\
8  & 70.1 & 66.4 \\
16 & 68.4 & \textbf{70.8} \\
32 & 66.7 & 72.1 \\
\bottomrule
\end{tabular}
\end{center}

Recurrent models peak at $D=4$; non-recurrent models keep improving
through $D=32$. At the recurrent optimum ($D=4$), the recurrent model
outperforms the non-recurrent model by 12.6 points at the same FLOPs.

\subsection*{FineWeb Language Modeling (Validation Perplexity, $\downarrow$)}

\begin{center}
\begin{tabular}{lcc}
\toprule
Depth $D$ & Recurrent & Non-Recurrent \\
\midrule
2  & 18.7 & 22.4 \\
4  & \textbf{17.3} & 20.1 \\
8  & 17.8 & 18.6 \\
16 & 18.4 & \textbf{17.9} \\
32 & 19.1 & 17.2 \\
\bottomrule
\end{tabular}
\end{center}

Consistent finding: recurrent optimal at $D=4$; non-recurrent at
$D=16$--$32$.

\section{Ablations and Interpretation}

\textbf{Width compensation:} Reducing recurrent depth from 8 to 2
layers while increasing hidden dimension to match FLOPs recovers 97\%
of 8-layer recurrent performance on both tasks. Width is a suitable
substitute for depth in recurrent networks.

\textbf{Expert count:} Increasing parallel experts from 1 to 4 (same
compute via width reduction) provides consistent modest gains for both
architectures, orthogonal to the depth effect.

\textbf{Shared vs.\ unshared weights:} Unshared weights slightly
improve non-recurrent baselines but do not change the qualitative
finding for recurrent models.

\textbf{Sequence length effect:} Longer sequences amplify the
advantage of shallow recurrent models over deep ones, consistent with
the interpretation that temporal depth substitutes for within-step
depth; more recurrent steps make additional layers more redundant.

\section*{Reference}

Ivan Anokhin, Johan Obando-Ceron, Irina Rish, Sebastian Risi.
\textbf{Temporal Recurrence Favors Fewer Layers.}
arXiv:2609.12531, September 2026.
\url{https://arxiv.org/abs/2609.12531}

\end{document}
